\section{Columnar storage}
\label{sec:pg-storage-columnar}

ClickBench is a benchmark for analytical DBMS and it also reported results for both Hydra and Citus. Figure~\ref{fig:storage-clickbench-motivation} uses historical benchmark submissions on AWS machines~\autocite{clickbench_2023_columnar}. 
The benchmark consisted of around 43 queries over a dataset of around 100 million rows (3 runs with minum of second two taken as into the sum).

\begin{figure}[htbp]
\centering
\begin{tikzpicture}[x=1cm,y=1cm,font=\footnotesize]
  \node at (5.2,4.0) {Hot-query time (min)};
  \node at (11.2,4.0) {Stored size (GB)};
  \foreach \v in {0,40,80,120,160} {
    \draw[black!12] (3+.0275*\v,.1)--(3+.0275*\v,3.55);
    \node[below] at (3+.0275*\v,.05) {\v};
  }
  \foreach \v in {0,20,40,60,80} {
    \draw[black!12] (9+.055*\v,.1)--(9+.055*\v,3.55);
    \node[below] at (9+.055*\v,.05) {\v};
  }
  % From the pinned ClickBench result files: sum(min(run2,run3))/60;
  % storage bytes / 1e9. All 43 queries are included, without penalties.
  \foreach \y/\t/\d/\tl/\dl/\col in {
    3.1/157.0344/77.797067741/157.0/77.8/storageBlue,
    1.9/38.596661767/18.980918899/38.6/19.0/storageOrange,
    .7/7.983093133/18.995669982/8.0/19.0/storageGreen} {
    \draw[draw=\col,fill=\col!25]
      (3,\y-.25) rectangle (3+.0275*\t,\y+.25);
    \node[anchor=west] at (3+.0275*\t+.08,\y) {\tl};
    \draw[draw=\col,fill=\col!25]
      (9,\y-.25) rectangle (9+.055*\d,\y+.25);
    \node[anchor=west] at (9+.055*\d+.08,\y) {\dl};
  }
  \node[anchor=east,align=right] at (2.75,3.1) {PostgreSQL 14\\2022 submission};
  \node[anchor=east,align=right] at (2.75,1.9) {Citus 11.0\\2022 submission};
  \node[anchor=east,align=right] at (2.75,.7) {Hydra Columnar\\2023 submission};
  \draw[black!55] (3,.1)--(7.4,.1);
  \draw[black!55] (9,.1)--(13.4,.1);
  \node at (5.2,-.65) {Sum over all 43 queries};
  \node at (11.2,-.65) {$1\ \mathrm{GB}=10^9$ bytes};
\end{tikzpicture}
\caption{Historical ClickBench, calculated from the
repository snapshot of 20 September 2023. Each submission uses
\texttt{c6a.4xlarge} machine with 16 vCPUs, 32~GiB RAM, and 500~GB of GP2
storage. PostgreSQL and Citus
results are dated 1 July 2022, and Hydra results 19 September 2023.
The times use the sum of shorter of runs two and three for every query. Note that these are "best-fit" results so the configuration was different. It's good however for quick comparison of performance.}
\label{fig:storage-clickbench-motivation}
\end{figure}

Both columnar submissions store about 19~GB, compared with PostgreSQL's
77.8~GB. This was due to support of compression on stored columns. Hydra~\autocite{hydra_columnar_source} had much better performance due to introduction of vectorized SQL execution for aggregates and column cache that wasn't present in Citus~\autocite{citus_distributed_2021}.
These results do not translate directly to BioseqDB as the benchmarks tests some arbitrarily-built queries, however it motivated me to put big effort into modernizing Hydra codebase and incorporating it into BioseqDB as mentioned in previous section~\ref{sec:pg-storage-role}.

\subsection{Columnar layout}


% %Stripe size is not configurable and is set to size of disk block minus page header so it can exactly fit into one disk block (around $8KB$, but depends on compile-time value Postgres version provides).
% \item A \textbf{stripe} is a batch of appended values, where values can be anything. In case of the user table these are grouped metadata values (native Postgres types) and for sequence columns - sequence locators. For internal sharded storage tables these are sequence fragments itself.
%     It's important to note that TOAST can have a limit up to $1GB$ per row, so large values (not only long sequence but also if user stores large BLOB) can span multiple stripes.
  

The columnar layout has three levels:

\begin{itemize}
  \item A \textbf{stripe} is a batch of chunks coresponding to set of rows. Normally one stripe holds maximum is 150,000 rows.
  \item A \textbf{chunk group} is a group of chunks for the same rows inside a stripe. By default one chunk group holds maximum of 10,000 rows.
  \item A \textbf{column chunk} is the most basic entity that just holds values for one group as a one column. It also stores aggregated values (count, min, max) to enable 
  faster aggregate operations (vectorization) and the bitmask (marking which values are NULLs).
\end{itemize}

Stripe and chunk-group sizes are options per table, where larger ones allows more efficient compression and smaller ones allows more fine-grained read access~\autocite{citus_columnar_compression_2021}.

Figure~\ref{fig:storage-heap-to-stripe} shows this model for the six rows.
By default Postgres row-based storage (heap tuple) stores \texttt{id}, \texttt{sample\_id}, and the
sequence locator together. The stripe is columnar representation of these rows.
Each column occupies its own sequence of chunks. A chunk group is only an abstract concept for chunks that are related to the same rows.
The last group of a stripe can be shorter if total row count isn't a multiple of chunk group size.

\begin{figure}[H]
\centering
\begin{tikzpicture}[x=1cm,y=1cm,font=\footnotesize]
  \node[anchor=west,font=\bfseries] at (0,6.10)
    {Heap tuples become a stripe of column chunks};
  \node[anchor=north west,align=left] at (0,5.75)
    {Heap: one tuple stores\\every column};
  \foreach \y/\id/\sid/\rn in {
    4.45/1/12/1,
    3.80/2/18/2,
    3.15/3/7/3,
    2.50/4/41/4,
    1.85/5/9/5,
    1.20/6/22/6} {
    \node[anchor=east] at (1.05,\y) {row \rn};
    \node[storage cell,minimum width=.62cm,fill=storageBlue!15]
      at (1.55,\y) {\id};
    \node[storage cell,minimum width=.62cm,fill=storageOrange!15]
      at (2.28,\y) {\sid};
    \node[storage cell,minimum width=.78cm,fill=storageGreen!15]
      at (3.15,\y) {loc};
  }
  \draw[storage arrow] (3.90,2.85)--(4.70,2.85);
  \node[align=center] at (4.30,3.35) {same\\rows};
  \draw[draw=black!45] (4.85,0.48) rectangle (12.15,5.48);
  \node[anchor=north,align=center,font=\bfseries] at (8.50,5.32)
    {One stripe\\every column of these rows};
  \node at (7.40,4.12) {\texttt{id}};
  \node at (9.20,4.12) {\texttt{sample\_id}};
  \node at (11.00,4.12) {\texttt{seq}};
  \node[anchor=west] at (5.05,3.40) {group 0};
  \node[anchor=west,text=storageOrange,font=\scriptsize] at (5.05,3.08)
    {deleted $0$};
  \node[anchor=west] at (5.05,1.38) {group 1};
  \node[anchor=west,text=storageOrange,font=\scriptsize] at (5.05,1.06)
    {deleted $0$};
  \node[storage cell,text width=1.45cm,inner sep=2pt,
    minimum height=.90cm,align=center,
    fill=storageBlue!15,draw=storageBlue,line width=1.1pt]
    (idchunk) at (7.40,3.25) {1, 2, 3};
  \node[storage cell,text width=1.45cm,inner sep=2pt,
    minimum height=.90cm,align=center,fill=storageOrange!15]
    at (9.20,3.25) {12, 18, 7};
  \node[storage cell,text width=1.45cm,inner sep=2pt,
    minimum height=.90cm,align=center,fill=storageGreen!15]
    at (11.00,3.25) {\scriptsize loc, loc, loc};
  \node[storage cell,text width=1.45cm,inner sep=2pt,
    minimum height=.90cm,align=center,fill=storageBlue!15]
    at (7.40,1.23) {4, 5, 6};
  \node[storage cell,text width=1.45cm,inner sep=2pt,
    minimum height=.90cm,align=center,fill=storageOrange!15]
    at (9.20,1.23) {41, 9, 22};
  \node[storage cell,text width=1.45cm,inner sep=2pt,
    minimum height=.90cm,align=center,fill=storageGreen!15]
    at (11.00,1.23) {\scriptsize loc, loc, loc};
  \node[draw=storageBlue,fill=storageBlue!6,minimum width=12.00cm,
    minimum height=1.62cm,inner sep=0pt,anchor=north west]
    (chunkdetail) at (0.15,0.10) {};
  \node[anchor=west,font=\bfseries,text=storageBlue] at (0.35,-0.08)
    {One column chunk};
  \node[anchor=west,text=black!65] at (4.55,-0.08)
    {\texttt{id} of group 0};
  \node at (2.25,-0.48) {presence bits};
  \foreach \i/\b in {0/1,1/1,2/1}
    \node[storage cell,minimum width=.50cm,inner sep=2pt,fill=storageBlue!12]
      at (1.55+0.70*\i,-1.08) {\b};
  \node at (5.70,-0.48) {value buffer};
  \foreach \i/\v in {0/1,1/2,2/3}
    \node[storage cell,minimum width=.50cm,inner sep=2pt,fill=storageGreen!13]
      at (5.00+0.70*\i,-1.08) {\v};
  \node at (9.70,-0.48) {aggregated values};
  \node[storage cell,minimum width=2.35cm,inner sep=1pt,align=center,
    fill=storageOrange!12] at (9.70,-1.08)
    {min $1$, max $3$\\count $3$};
  \node[draw=storageOrange,fill=storageOrange!8,minimum width=12.00cm,
    minimum height=1.28cm,inner sep=0pt,anchor=north west]
    (groupdetail) at (0.15,-1.68) {};
  \node[anchor=west,font=\bfseries,text=storageOrange] at (0.35,-1.88)
    {One chunk group};
  \node[anchor=west,text=black!65] at (4.55,-1.88)
    {shared by every column};
  \node at (2.25,-2.28) {delete bits};
  \foreach \i/\b in {0/0,1/0,2/0}
    \node[storage cell,minimum width=.50cm,inner sep=2pt,fill=storageOrange!12]
      at (1.55+0.70*\i,-2.68) {\b};
  \node at (7.20,-2.28) {deleted rows};
  \node[storage cell,minimum width=1.20cm,inner sep=2pt,fill=storageOrange!12]
    at (7.20,-2.68) {$0$};
  \draw[storage arrow,storageBlue]
    (idchunk.south) -- (7.40,2.18) -- (12.60,2.18) -- (12.60,-0.71)
    -- (chunkdetail.east);
  \node[fill=white,inner sep=2pt,text=storageBlue] at (9.90,2.18)
    {One column chunk};
\end{tikzpicture}
\caption{The same six rows as heap tuples and as one stripe. Matching
colours mark the same column. The blue callout is one column chunk: presence
bits, present values, min/max, and slot count. The orange callout is the
chunk group: a deletion mask and \texttt{deleted\_rows}}
\label{fig:storage-heap-to-stripe}
\end{figure}


When I want to store row, each time the writer append this row to the current chunk group (in memory).
After sufficient number of row is reached (max chunk group size) the code serializes each column chunk to a stripe buffer that still is in memory
If the number of rows per stripe is reached (max stripe size) then stripe is "flushed", which means writing the stripe.
Commonly the maximums of stripe row counts will be hit, but there's additional check if the current chunk group for specific column buffer exceeds $1GB$ limit, then it forces flush of the entire stripe.
This means that if we store large BLOB along very small column, the BLOB will force the stripe to end i.e stripes always need to account for the same number of rows in each column (this also implies that the decision to store sequences via locators increases efficiency of that storage model as more tuples with non-sequence fields can be packed into single stripe).
Additionally there's hard limit on maximum size of $256MiB$ for a single value in row.

\warn{Storing large BLOBs}{
  Users should avoid storing large BLOBs along small columns - either store BLOBs of similar sizes together or only store smaller payloads in the same table.
  This is potential optimisation to be implemented in the future, although it's very low priority.
}

The figure~\ref{fig:storage-write-path} shows 3 main stages: we save tuples first to accumulate enough of them, then serialize them to a buffer and when stripe's buffer eaches $1GB$ or the strip max row count, only then it's flushed.
Algorithm~\ref{alg:storage-flush-stripe} presents pseudocode for that process as well.

\begin{figure}[H]
\centering
\begin{tikzpicture}[x=1cm,y=1cm,font=\footnotesize]
  \node[anchor=west,font=\bfseries] at (0,3.05)
    {Write path: row, then group, then stripe};
  \node[draw=black!45,fill=storageBlue!6,minimum width=3.55cm,
    minimum height=2.42cm,inner sep=0pt,anchor=north west]
    (stage1) at (0,2.70) {};
  \node[anchor=north,align=center,font=\bfseries,text=storageBlue]
    at (1.78,2.55) {1. Current group};
  \node[anchor=north,font=\scriptsize,text=black!60] at (1.78,2.22) {RAM};
  \node[storage cell,minimum width=.62cm,inner sep=1pt,fill=storageBlue!15]
    at (0.68,1.45) {id};
  \node[storage cell,minimum width=.70cm,inner sep=1pt,fill=storageOrange!15]
    at (1.78,1.45) {sid};
  \node[storage cell,minimum width=.62cm,inner sep=1pt,fill=storageGreen!15]
    at (2.88,1.45) {loc};
  \node[align=center,font=\scriptsize] at (1.78,0.72)
    {one INSERT writes\\one slot in every column};
  \draw[storage arrow] (3.70,1.45)--(4.35,1.45);
  \node[align=center,font=\scriptsize] at (4.02,1.95) {group\\full};
  \node[draw=black!45,fill=storageOrange!6,minimum width=3.70cm,
    minimum height=2.42cm,inner sep=0pt,anchor=north west]
    (stage2) at (4.50,2.70) {};
  \node[anchor=north,align=center,font=\bfseries,text=storageOrange]
    at (6.35,2.55) {2. Stripe buffers};
  \node[anchor=north,font=\scriptsize,text=black!60] at (6.35,2.22) {still RAM};
  \foreach \x/\fill/\lab in {
    5.45/storageBlue!15/id,
    6.50/storageOrange!15/sid,
    7.55/storageGreen!15/loc} {
    \node[storage cell,minimum width=.72cm,minimum height=.42cm,
      inner sep=1pt,fill=\fill] at (\x,1.68) {\scriptsize\lab};
    \node[storage cell,minimum width=.72cm,minimum height=.42cm,
      inner sep=1pt,fill=\fill] at (\x,1.12) {};
  }
  \node[font=\scriptsize,anchor=west] at (4.62,1.68) {g0};
  \node[font=\scriptsize,anchor=west] at (4.62,1.12) {g1};
  \node[align=center,font=\scriptsize] at (6.35,0.62)
    {pack each column,\\then compress};
  \draw[storage arrow] (8.35,1.45)--(9.00,1.45);
  \node[align=center,font=\scriptsize] at (8.68,2.05)
    {$S$ rows,\\$1$\,GB,\\or close};
  \node[draw=black!45,fill=storageGreen!6,minimum width=3.20cm,
    minimum height=2.42cm,inner sep=0pt,anchor=north west]
    (stage3) at (9.15,2.70) {};
  \node[anchor=north,align=center,font=\bfseries,text=storageGreen]
    at (10.75,2.55) {3. Flush};
  \node[anchor=north,font=\scriptsize,text=black!60] at (10.75,2.22) {disk};
  \node[draw=black!50,fill=white,minimum width=1.55cm,minimum height=.42cm,
    inner sep=0pt] at (10.20,1.68) {};
  \node[draw=black!50,fill=white,minimum width=1.55cm,minimum height=.42cm,
    inner sep=0pt] at (10.20,1.18) {};
  \node[font=\scriptsize] at (10.20,1.68) {pages};
  \node[font=\scriptsize] at (10.20,1.18) {catalogs};
  \node[align=center,font=\scriptsize] at (10.75,0.62)
    {all groups of\\this stripe};
\end{tikzpicture}
\caption{How we pack chunks groups into stripe buffers, then to flush them if necessary.}
\label{fig:storage-write-path}
\end{figure}

\begingroup
\captionsetup{type=algorithm}
\captionof{algorithm}{Appending rows}
\label{alg:storage-flush-stripe}
\small
\begin{algorithmic}[1]
  \Require columns $C$, group limit $G$ (default $10{,}000$), stripe limit $S$
    (default $150{,}000$), comparators $\mathrm{Compare}_c$ (possibly
    $\mathrm{NULL}$)
  \State $G_{\max}\gets\lceil S/G\rceil$
  \State $s\gets\mathrm{NULL}$ \Comment{open stripe, or none}
  \Statex
  \Procedure{AllocateNewStripe}{}
    \State $s.n\gets 0$ \Comment{rows already in this stripe}
    \State $s.P\gets 0^{|C|\times G}$
      \Comment{presence bits of the current group}
    \State $s.B\gets\varepsilon^{|C|}$
      \Comment{serialized present values of the current group}
    \State $s.K\gets\{\min:\bot,\max:\bot,\mathit{cnt}:0\}^{|C|\times G_{\max}}$
      \Comment{skip min, max, and slot count of every column chunk}
    \State $s.\mathrm{chunk}\gets\{P:\varepsilon,V:\varepsilon\}^{|C|\times G_{\max}}$
      \Comment{packed streams of every serialized group}
    \State $s.N\gets 0^{G_{\max}}$
      \Comment{row count of each serialized group}
    \State \Return $s$
  \EndProcedure
  \Statex
  \Procedure{WriteRow}{$v$}
    \If{$\mathrm{IsOpen}(s)\land\exists c\in C:\mathrm{Size}(s.B_c\Vert v_c)>1\,\mathrm{GB}$}
      \State \Call{FlushStripe}{$s$}
    \EndIf
    \If{$\neg\mathrm{IsOpen}(s)$}
      \State $s\gets\Call{AllocateNewStripe}{}$
    \EndIf
    \State $g\gets\lfloor s.n/G\rfloor$,\quad $r\gets s.n\bmod G$
    \ForAll{$c\in C$}
      \If{$\mathrm{IsNull}(v_c)$}
        \State $s.P_{c,r}\gets 0$
      \Else
        \If{$\mathrm{Align}(v_c)\ge 256\,\mathrm{MiB}$}
          \State \textbf{error}
        \EndIf
        \State $s.P_{c,r}\gets 1$
        \State $s.B_c\gets s.B_c\Vert\mathrm{Serialize}(v_c)$
        \If{$\mathrm{Compare}_c\neq\mathrm{NULL}$}
          \State $s.K_{c,g}.\min\gets\min_{\mathrm{Compare}_c}(s.K_{c,g}.\min,v_c)$
          \State $s.K_{c,g}.\max\gets\max_{\mathrm{Compare}_c}(s.K_{c,g}.\max,v_c)$
            \Comment{$\bot$ is the identity}
        \EndIf
      \EndIf
      \State $s.K_{c,g}.\mathit{cnt}\gets s.K_{c,g}.\mathit{cnt}+1$
    \EndFor
    \If{$r=G-1$}
      \State \Call{SerializeGroup}{$s$, $g$, $G$}
    \EndIf
    \State $s.n\gets s.n+1$
    \If{$s.n\ge S$}
      \State \Call{FlushStripe}{$s$}
    \EndIf
  \EndProcedure
  \Statex
  \Procedure{SerializeGroup}{$s$, $g$, $k$}
    \State $s.N_g\gets k$
    \ForAll{$c\in C$}
      \State $s.\mathrm{chunk}_{c,g}.P\gets
        \mathrm{PackBits}(\langle s.P_{c,r}\rangle_{r=0}^{k-1})$
    \EndFor
    \ForAll{$c\in C$}
      \State $s.\mathrm{chunk}_{c,g}.V\gets\mathrm{Compress}(s.B_c)$
        \Comment{identity if incompressible}
      \State $s.B_c\gets\varepsilon$
    \EndFor
  \EndProcedure
  \Statex
  \Procedure{FlushStripe}{$s$}
    \State $k\gets s.n\bmod G$
    \If{$k>0$}
      \State \Call{SerializeGroup}{$s$, $\lfloor s.n/G\rfloor$, $k$}
    \EndIf
    \State $m\gets\lceil s.n/G\rceil$
    \ForAll{$c\in C$}
      \State $\mathrm{Write}(\langle s.\mathrm{chunk}_{c,g}.P\rangle_{g=0}^{m-1})$
      \State $\mathrm{Write}(\langle s.\mathrm{chunk}_{c,g}.V\rangle_{g=0}^{m-1})$
    \EndFor
    \State \Call{WriteCatalog}{$s$}
    \State $s\gets\mathrm{NULL}$
  \EndProcedure
  \Statex
  \Procedure{WriteCatalog}{$s$}
    \State $\mathrm{Write}(\langle s.N, s.K, 0^{s.n} \rangle)$
      \Comment{Saves row count for group, aggegate stats and deletion mask}
  \EndProcedure
\end{algorithmic}
\endgroup

In's important that closing the writer also flushes any leftover stripe, so all in-memory buffers are saved to disk when transaction is committed.

As presented on figure~\ref{fig:storage-heap-to-stripe} column chunk stores a presence bitmap, a value buffer of only the
present values, and aggregated values used later for delettion and vectorized aggregation: the smallest
and largest non-NULL values (when the type has a comparison function), and a
row count (\texttt{value\_count}) counting all values in the chunk including NULLs.

To support deletion, the chunk group stores \texttt{deleted\_rows}, along with separate row mask representing which rows were deleted.
\texttt{DELETE} only sets the matching mask bit in the chunk group's mask and increments \texttt{deleted\_rows}. This allows BioseqDB to offer fast deletion without the need 
to rewrite the columns themselves. The wasted storage can be however reclaimed:
\begin{itemize}
  \item \texttt{VACUUM (TRUNCATE)} may rewrite last stripes in table whose undeleted rows still fit in
    one stripe and whose decompressed size is below $1\,\mathrm{GB}$. A single
    last stripe is rewritten only when more than $20\%$ of its rows are
    deleted. The rewrite copies undeleted rows into a new stripe and drops the old ones.
  \item The internal tables storing sequence fragments cannot be rewritten as each locator needs to store the row number. Supporting VACUUM for sequence fragments would need to be dispatched to all tables including user table. Currently VACUUM is only supported on the user table.
  \item \texttt{VACUUM FULL} can rewrite the whole table to reclaim wasted space.
\end{itemize}

\subsection{Smart reads}

In BioseqDB it's unnecessary to read all chunks all the time. The two optimisations are:
\begin{enumerate}
  \item \emph{Projection} means that only selected columns from \texttt{SELECT} query are red. This is trivial optimisation.
  \item \emph{Min/max pruning} is more interesting one. For an ordered column,
    the writer records the smallest and largest non-NULL values in each chunk (in aggregates fields seen in figure~\ref{fig:storage-heap-to-stripe}).
    The BioseqDB TAM presented in figure~\ref{fig:bioseqscan-hook-lifecycle} can determine if none of the values in the chunk group satisfy the filter.
    If that happens, entire chunk group is skipped.
    For example, a group with sample identifiers from 1 to 100 cannot satisfy \texttt{BETWEEN 500 AND 600}.
\label{lst:storage-minmax-pruning}
\end{enumerate}

Figure~\ref{fig:column-store-layout} shows both optimisations. Group that satisfies both of these conditions still needs to apply filters to each row individually, but I save time but not touching groups that are not revelant to the search.
The figure assumes query \texttt{SELECT id FROM table WHERE sample\_id BETWEEN 500 AND 600} is executed.

\begin{figure}[H]
\centering
\begin{tikzpicture}[x=1cm,y=1cm,font=\footnotesize]
  \node[anchor=west,font=\bfseries] at (0,4.45)
    {One stripe: three chunk groups and three columns};
  \node at (3.3,3.85) {\texttt{id}};
  \node at (6.15,3.85) {\texttt{sample\_id}};
  \node at (9.0,3.85) {\texttt{seq}};
  \foreach \y/\g in {3.05/0,1.95/1,.85/2} {
    \node[anchor=east] at (1.7,\y) {group \g};
    \foreach \x in {1.9,4.75,7.6}
      \draw[draw=black!45,fill=black!5] (\x,\y-.4)
        rectangle +(2.65,.8);
  }
  \draw[storageBlue,fill=storageBlue!15] (1.9,1.55) rectangle (4.55,2.35);
  \draw[storageBlue,fill=storageBlue!15] (4.75,1.55) rectangle (7.4,2.35);
  \node at (3.23,1.95) {read identifiers};
  \node at (6.08,1.95) {bounds $[450,700]$};
  \node at (6.08,3.05) {bounds $[1,100]$};
  \node at (6.08,.85) {bounds $[800,950]$};
  \foreach \y in {3.05,.85} {
    \draw[storageOrange,line width=1.1pt] (2.05,\y-.28)--(4.4,\y+.28);
    \draw[storageOrange,line width=1.1pt] (2.05,\y+.28)--(4.4,\y-.28);
    \node[anchor=west,text=storageOrange] at (10.55,\y) {skip group};
  }
  \node[align=center,text=black!65] at (8.93,1.95) {locator stream\\not requested};
  \node[anchor=west,text=storageBlue,align=left] at (10.55,1.95)
    {read group;\\test its rows};
  \draw[decorate,decoration={brace,mirror,amplitude=4pt},black!60]
    (7.6,.22)--(10.25,.22);
  \node[align=center] at (8.93,-.37) {no sequence\\fragment reads};
\end{tikzpicture}
\caption{Column selection and group pruning for a simple query \texttt{SELECT id FROM table WHERE sample\_id BETWEEN 500 AND 600}.
Blue cells are read, orange crosses mark groups whose aggregates exclude all the rows. Unrequested sequence locators are unread,
because they were not specified in the \texttt{SELECT} clause.
In the drawing I only illustrate the idea of reading rows, not their actual layout in storage.}
\label{fig:column-store-layout}
\end{figure}

\info{Prunning for sequence columns}{
  Sequence columns with locators are excluded for min/max pruning, but it's possible to implement that feature theoretically. 
  I decided not to do it, because storing min/max values would be simple as locator isn't big object, however computing these values for very large sequences
  will be computionally expensive and would negatively impact table inserts.
}

\subsection{Compression}

Figure~\ref{fig:storage-column-streams} shows in detail how presence bitmap is computed.
If the NULLs are very frequent, I will save the space by not storing them.

\begin{figure}[H]
\centering
\begin{tikzpicture}[x=1cm,y=1cm,font=\footnotesize]
  \node[anchor=east] at (2.1,3.15) {row position};
  \node[anchor=east] at (2.1,2.35) {column value};
  \node[anchor=east] at (2.1,1.55) {presence bit};
  \foreach \i/\v/\b in {0/12/1,1/NULL/0,2/17/1,3/17/1,4/NULL/0,5/21/1} {
    \node at (2.75+1.35*\i,3.15) {\i};
    \node[storage cell,minimum width=1.2cm,fill=black!4]
      at (2.75+1.35*\i,2.35) {\v};
    \node[storage cell,minimum width=1.2cm,fill=storageBlue!12]
      at (2.75+1.35*\i,1.55) {\b};
  }
  \node[anchor=east] at (2.1,.25) {value buffer};
  \foreach \i/\v in {0/12,1/17,2/17,3/21}
    \node[storage cell,minimum width=1.2cm,fill=storageGreen!13]
      at (2.75+1.35*\i,.25) {\v};
  \draw[storage arrow] (8.0,.25)--(9.0,.25);
  \node[storage cell,fill=storageOrange!12,minimum width=2.25cm]
    at (10.25,.25) {compress values};
  \node[anchor=west,text=black!65] at (2.15,-.5)
    {Only four values are stored; the bitmap keeps their six row positions.};
\end{tikzpicture}
\caption{Presence bitmap is useful to enable better NULL compression}
\label{fig:storage-column-streams}
\end{figure}

The Hydra code that BioseqDB columnar TAM was based on supported PGLZ, LZ4, or Zstandard, depending on the
build and table options. In BioseqDB all compression is disabled by default as I would like to benchmark in the future best (and more modern) compression algorithms decicated for example to sequence storage.
Abadi et al. show that the choice of compression method depends on both
the data and the query workload~\autocite{abadi2006compression}, so I leave it as future extension.

\section{Raw byte storage}
\label{sec:pg-storage-pages}

After stripe is serialized and flushed it needs to be written to disk - function $Write$ and $WriteCatalog$ in the pseudocode~\ref{alg:storage-flush-stripe}
\begin{itemize}
  \item Physically first written is the presence bitmap, then the value buffer for a column, then each column follow. These are placed in the user relation database pages.
  \item Catalog (row counts, aggregates and deletion mask) is written separately from the stripe data and it's placed in usual heap rows i.e BioseqDB introduces separate relations for storing chunks information.
\end{itemize}

Catalog is represented by four relations inside the BioseqDB:

\begin{minipage}{\linewidth}
\begin{lstlisting}[style=bioseq-storage-sql,caption={Metadata catalogs for flushed stripes.},label={lst:storage-metadata-catalogs}]
CREATE TABLE bioseq_internal.stripe (
  storage_id bigint,
  stripe_num bigint,
  file_offset bigint,
  data_length bigint,
  column_count integer,
  chunk_row_count integer,
  row_count bigint,
  chunk_group_count integer,
  first_row_number bigint
);

CREATE TABLE bioseq_internal.chunk_group (
  storage_id bigint,
  stripe_num bigint,
  chunk_group_num integer,
  row_count bigint,
  deleted_rows bigint
);

CREATE TABLE bioseq_internal.chunk (
  storage_id bigint,
  stripe_num bigint,
  attr_num integer,
  chunk_group_num integer,
  minimum_value bytea,
  maximum_value bytea,
  value_stream_offset bigint,
  value_stream_length bigint,
  exists_stream_offset bigint,
  exists_stream_length bigint,
  value_compression_type integer,
  value_compression_level integer,
  value_decompressed_length bigint,
  value_count bigint
);

CREATE TABLE bioseq_internal.row_mask (
  id bigint,
  storage_id bigint,
  stripe_id bigint,
  chunk_id integer,
  start_row_number bigint,
  end_row_number bigint,
  deleted_rows integer,
  mask bytea
);
\end{lstlisting}
\end{minipage}

When writing stripe data I need to take care about page boundaries.
Let $\mathrm{BLCKSZ}$ be the PostgreSQL page size in bytes (this is compile-time constant provided by Postgres), $H$ the
page-header size (technically $\mathrm{sizeof(PageHeaderData)}$)~\autocite{postgresql17_storage_page_layout}, and $P=\mathrm{BLCKSZ}-H$ the remaining page capacity.
If $x\geq0$ is the total offset of the $x$-th byte, then page coordinates are defined as:
\[
  \operatorname{block}(x)=2+\left\lfloor x/P\right\rfloor,
  \qquad
  \operatorname{position}(x)=H+(x\bmod P).
\]
BioseqDB will write multiple pages for each stripe, so I need to take additional care about buffers going through multiple page boundaries as seen on figure~\ref{fig:storage-pages}.

\begin{figure}[htbp]
\centering
\begin{tikzpicture}[x=1cm,y=1cm,font=\footnotesize]
  \node[anchor=west,font=\bfseries] at (0,4.05)
    {Toy pages: 2 header bytes and 10 payload bytes};
  \node[anchor=east] at (1.45,2.85) {logical bytes};
  \draw[draw=black!50,fill=black!5] (1.75,2.5) rectangle (11.75,3.2);
  \draw[storageBlue,fill=storageBlue!20] (5.75,2.5) rectangle (9.25,3.2);
  \foreach \x/\v in {1.75/0,5.75/8,6.75/10,9.25/15,11.75/20}
    \node[above=2pt] at (\x,3.2) {$\v$};
  \node at (7.5,2.85) {next buffer fragment $[8,15)$ coming from $\mathrm{Write}(\langle s.\mathrm{chunk}_{c,g}.P\rangle_{g=0}^{m-1}$};
  \node[anchor=east] at (1.45,1.15) {stored pages};
  \draw[fill=black!20,draw=black!60] (1.75,.8) rectangle (2.55,1.5);
  \draw[fill=black!5,draw=black!60] (2.55,.8) rectangle (6.55,1.5);
  \draw[storageBlue,fill=storageBlue!20] (5.75,.8) rectangle (6.55,1.5);
  \draw[fill=black!20,draw=black!60] (7.0,.8) rectangle (7.8,1.5);
  \draw[fill=black!5,draw=black!60] (7.8,.8) rectangle (11.8,1.5);
  \draw[storageBlue,fill=storageBlue!20] (7.8,.8) rectangle (9.8,1.5);
  \node at (2.15,1.15) {H};
  \node at (7.4,1.15) {H};
  \node at (6.15,1.15) {2};
  \node at (8.8,1.15) {5 bytes};
  \node at (4.15,.35) {block 2};
  \node at (9.4,.35) {block 3};
  \draw[storage arrow,storageBlue] (6.25,2.45)--(6.15,1.55);
  \draw[storage arrow,storageBlue] (8.0,2.45)--(8.8,1.55);
  \node[anchor=west,text=black!65] at (1.75,-.3)
    {Copy two bytes, skip the next page header, then copy five bytes.};
\end{tikzpicture}
\caption{I use logical offsets $\operatorname{block}(x)$ and $\operatorname{position}(x)$ to hide page boundaries. Diagram shows that we want to write some fragment
of in-memory buffer $[8,15)$ to the disk, but it needs to be split across multiple pages. In reality the dimensions are larger than just 20 bytes, but it illustrates the idea.}
\label{fig:storage-pages}
\end{figure}

The original idea comes from Citus page writer that it's avaialble in their public repository\coderef
{citusstorage/src/backend/columnar/columnar_storage.c}
{Citus columnar page mapping}{columnar_storage.c}{118}

To allow concurrent access (queries during inserts) I need to note that the writer does the following things:
\begin{itemize}
  \item \textbf{Reservation:} Identifiers for chunks and groups is done under the relation extension lock.
  \item \textbf{Write:} Writes to pages are done using PostgreSQL buffer locks and generic write-ahead logging (WAL)~\autocite{postgresql17_generic_wal}.
  \item \textbf{Record completed data:} The logic uses Postgres' Snapshots so data will be visible only after the transaction is committed.
\end{itemize}
